\chapter{控制器设计}\label{chap:controller}

本章在第三章建模结果的基础上给出整机控制器设计。双轮足机器人需要同时满足矢状面内的倒立摆平衡与行进控制、腿部支撑与姿态调节，以及侧向滚转稳定等目标。为兼顾实时性与可调性，本文采用分层混合控制结构：矢状面主通道基于线性化状态空间模型设计 LQR 平衡控制器，输出轮系驱动力矩与腿部公共摆动量；腿部通道在任务空间内采用 VMC，分别调节腿长与摆角；侧向通道引入滚转补偿，对左右腿支撑力做差动修正。各通道控制量在执行层融合后下发至驱动器。

\section{线性化状态空间方程推导}

\subsection{广义坐标与输入变量选取}

为与前文整机建模结果保持一致，本文仍选取左右轮转角、左右腿摆角以及机体俯仰角作为独立广义坐标，即
\begin{equation}
  \boldsymbol{q} =
  \begin{bmatrix}
    \theta_{wl} & \theta_{wr} & \theta_{ll} & \theta_{lr} & \theta_b
  \end{bmatrix}^{\mathsf{T}},
  \label{eq:full-generalized-coordinate}
\end{equation}
并定义输入向量为
\begin{equation}
  \boldsymbol{u} =
  \begin{bmatrix}
    T_{lw,l} & T_{lw,r} & T_{bl,l} & T_{bl,r}
  \end{bmatrix}^{\mathsf{T}}.
  \label{eq:full-input-vector}
\end{equation}

需要指出的是，整机前进位移 $s$ 与偏航角 $\phi$ 并非独立广义坐标，而是由左右轮转角与腿部构型决定的几何派生量。本文先围绕 $\boldsymbol{q}$ 完成动力学线性化，再通过运动学映射把结果转换为更便于控制设计的物理状态变量形式。

\subsection{线性化广义动力学方程}

选取机器人静态直立工况作为平衡工作点，并在该点附近进行小角度近似与一阶扰动展开。忽略高阶小量后，第三章的非线性动力学方程可线性化为
\begin{equation}
  \boldsymbol{M}\ddot{\boldsymbol{q}} + \boldsymbol{K}\boldsymbol{q} = \boldsymbol{E}\boldsymbol{u},
  \label{eq:linearized-generalized-dynamics}
\end{equation}
其中，$\boldsymbol{M} \in \mathbb{R}^{5\times 5}$ 为平衡点处的等效惯性矩阵，$\boldsymbol{K} \in \mathbb{R}^{5\times 5}$ 为由重力项及约束耦合项线性化得到的刚度矩阵，$\boldsymbol{E} \in \mathbb{R}^{5\times 4}$ 为输入分配矩阵。

由于 $\boldsymbol{M}$ 在工作点邻域内非奇异，故式 \eqref{eq:linearized-generalized-dynamics} 可进一步改写为显式加速度形式
\begin{equation}
  \ddot{\boldsymbol{q}}
  =
  \boldsymbol{A}_{\mathrm{dyn}} \boldsymbol{q}
  +
  \boldsymbol{B}_{\mathrm{dyn}} \boldsymbol{u},
  \label{eq:qdd-dyn}
\end{equation}
其中
\begin{equation}
  \boldsymbol{A}_{\mathrm{dyn}} = -\boldsymbol{M}^{-1}\boldsymbol{K},
  \qquad
  \boldsymbol{B}_{\mathrm{dyn}} = \boldsymbol{M}^{-1}\boldsymbol{E}.
  \label{eq:Adyn-Bdyn}
\end{equation}

\subsection{运动学映射与状态变量构造}

根据第三章建立的几何关系，整机前进位移满足
\begin{equation}
  s = \frac{R_w}{2}\left(\theta_{wl} + \theta_{wr}\right),
  \qquad
  \ddot{s} = \frac{R_w}{2}\left(\ddot{\theta}_{wl} + \ddot{\theta}_{wr}\right).
  \label{eq:s-kinematic-map}
\end{equation}
偏航角可表示为
\begin{equation}
  \phi
  =
  \frac{R_w}{2R_l}\left(-\theta_{wl} + \theta_{wr}\right)
  -
  \frac{l_l}{2R_l}\sin \theta_{ll}
  +
  \frac{l_r}{2R_l}\sin \theta_{lr}.
  \label{eq:phi-kinematic-map}
\end{equation}
在平衡点附近进行小角度线性化后，可得
\begin{equation}
  \ddot{\phi}
  =
  -\frac{R_w}{2R_l}\ddot{\theta}_{wl}
  +
  \frac{R_w}{2R_l}\ddot{\theta}_{wr}
  -
  \frac{l_l}{2R_l}\ddot{\theta}_{ll}
  +
  \frac{l_r}{2R_l}\ddot{\theta}_{lr}.
  \label{eq:phi-acc-map}
\end{equation}

进一步定义加速度状态向量
\begin{equation}
  \boldsymbol{x}_{\mathrm{acc}}
  =
  \begin{bmatrix}
    \ddot{s} & \ddot{\phi} & \ddot{\theta}_{ll} & \ddot{\theta}_{lr} & \ddot{\theta}_b
  \end{bmatrix}^{\mathsf{T}},
  \label{eq:x-acc-def}
\end{equation}
则由式 \eqref{eq:s-kinematic-map} 和式 \eqref{eq:phi-acc-map} 可将其写成
\begin{equation}
  \boldsymbol{x}_{\mathrm{acc}}
  =
  \boldsymbol{T}_{\mathrm{map}}\ddot{\boldsymbol{q}},
  \label{eq:xacc-tmap}
\end{equation}
其中
\begin{equation}
  \boldsymbol{T}_{\mathrm{map}}
  =
  \begin{bmatrix}
    \dfrac{R_w}{2} & \dfrac{R_w}{2} & 0 & 0 & 0 \\
    -\dfrac{R_w}{2R_l} & \dfrac{R_w}{2R_l} & -\dfrac{l_l}{2R_l} & \dfrac{l_r}{2R_l} & 0 \\
    0 & 0 & 1 & 0 & 0 \\
    0 & 0 & 0 & 1 & 0 \\
    0 & 0 & 0 & 0 & 1
  \end{bmatrix}.
  \label{eq:tmap}
\end{equation}

为构造位置状态，再定义
\begin{equation}
  \boldsymbol{x}_{\mathrm{pos}}
  =
  \begin{bmatrix}
    s & \phi & \theta_{ll} & \theta_{lr} & \theta_b
  \end{bmatrix}^{\mathsf{T}}.
  \label{eq:xpos-def}
\end{equation}
由式 \eqref{eq:s-kinematic-map} 和式 \eqref{eq:phi-kinematic-map} 的线性化关系可反解得到
\begin{equation}
  \boldsymbol{q} = \boldsymbol{A}_{\mathrm{map}} \boldsymbol{x}_{\mathrm{pos}},
  \label{eq:q-amap-x}
\end{equation}
其中
\begin{equation}
  \boldsymbol{A}_{\mathrm{map}}
  =
  \begin{bmatrix}
    \dfrac{1}{R_w} & -\dfrac{R_l}{R_w} & -\dfrac{l_l}{2R_w} & \dfrac{l_r}{2R_w} & 0 \\
    \dfrac{1}{R_w} & \dfrac{R_l}{R_w} & \dfrac{l_l}{2R_w} & -\dfrac{l_r}{2R_w} & 0 \\
    0 & 0 & 1 & 0 & 0 \\
    0 & 0 & 0 & 1 & 0 \\
    0 & 0 & 0 & 0 & 1
  \end{bmatrix}.
  \label{eq:amap}
\end{equation}

将式 \eqref{eq:qdd-dyn}、式 \eqref{eq:xacc-tmap} 与式 \eqref{eq:q-amap-x} 联立，可得
\begin{equation}
  \boldsymbol{x}_{\mathrm{acc}}
  =
  \boldsymbol{T}_{\mathrm{map}}
  \left(
    \boldsymbol{A}_{\mathrm{dyn}}\boldsymbol{A}_{\mathrm{map}}\boldsymbol{x}_{\mathrm{pos}}
    +
    \boldsymbol{B}_{\mathrm{dyn}}\boldsymbol{u}
  \right).
  \label{eq:xacc-expand}
\end{equation}
定义
\begin{equation}
  \boldsymbol{A}_{\mathrm{eff}}
  =
  \boldsymbol{T}_{\mathrm{map}}\boldsymbol{A}_{\mathrm{dyn}}\boldsymbol{A}_{\mathrm{map}},
  \qquad
  \boldsymbol{B}_{\mathrm{eff}}
  =
  \boldsymbol{T}_{\mathrm{map}}\boldsymbol{B}_{\mathrm{dyn}},
  \label{eq:aeff-beff}
\end{equation}
则有
\begin{equation}
  \boldsymbol{x}_{\mathrm{acc}}
  =
  \boldsymbol{A}_{\mathrm{eff}}\boldsymbol{x}_{\mathrm{pos}}
  +
  \boldsymbol{B}_{\mathrm{eff}}\boldsymbol{u}.
  \label{eq:xacc-eff}
\end{equation}

\subsection{标准状态空间方程}

按本文所选状态排序方式，定义完整状态向量为
\begin{equation}
  \boldsymbol{x}
  =
  \begin{bmatrix}
    s & \dot{s} & \phi & \dot{\phi} & \theta_{ll} & \dot{\theta}_{ll} &
    \theta_{lr} & \dot{\theta}_{lr} & \theta_b & \dot{\theta}_b
  \end{bmatrix}^{\mathsf{T}}.
  \label{eq:full-state-vector}
\end{equation}
记 $\boldsymbol{A}_{\mathrm{eff}} = [a_{ij}] \in \mathbb{R}^{5\times 5}$，$\boldsymbol{B}_{\mathrm{eff}} = [b_{ij}] \in \mathbb{R}^{5\times 4}$，则线性化状态空间模型可写为
\begin{equation}
  \dot{\boldsymbol{x}} = \boldsymbol{A}\boldsymbol{x} + \boldsymbol{B}\boldsymbol{u},
  \label{eq:linear-state-space}
\end{equation}
其中
\begin{equation}
  \boldsymbol{A}
  =
  \begin{bmatrix}
    0 & 1 & 0 & 0 & 0 & 0 & 0 & 0 & 0 & 0 \\
    a_{11} & 0 & a_{12} & 0 & a_{13} & 0 & a_{14} & 0 & a_{15} & 0 \\
    0 & 0 & 0 & 1 & 0 & 0 & 0 & 0 & 0 & 0 \\
    a_{21} & 0 & a_{22} & 0 & a_{23} & 0 & a_{24} & 0 & a_{25} & 0 \\
    0 & 0 & 0 & 0 & 0 & 1 & 0 & 0 & 0 & 0 \\
    a_{31} & 0 & a_{32} & 0 & a_{33} & 0 & a_{34} & 0 & a_{35} & 0 \\
    0 & 0 & 0 & 0 & 0 & 0 & 0 & 1 & 0 & 0 \\
    a_{41} & 0 & a_{42} & 0 & a_{43} & 0 & a_{44} & 0 & a_{45} & 0 \\
    0 & 0 & 0 & 0 & 0 & 0 & 0 & 0 & 0 & 1 \\
    a_{51} & 0 & a_{52} & 0 & a_{53} & 0 & a_{54} & 0 & a_{55} & 0
  \end{bmatrix},
  \label{eq:state-matrix-A}
\end{equation}
\begin{equation}
  \boldsymbol{B}
  =
  \begin{bmatrix}
    0 & 0 & 0 & 0 \\
    b_{11} & b_{12} & b_{13} & b_{14} \\
    0 & 0 & 0 & 0 \\
    b_{21} & b_{22} & b_{23} & b_{24} \\
    0 & 0 & 0 & 0 \\
    b_{31} & b_{32} & b_{33} & b_{34} \\
    0 & 0 & 0 & 0 \\
    b_{41} & b_{42} & b_{43} & b_{44} \\
    0 & 0 & 0 & 0 \\
    b_{51} & b_{52} & b_{53} & b_{54}
  \end{bmatrix}.
  \label{eq:state-matrix-B}
\end{equation}

式 \eqref{eq:linear-state-space} 即为本文用于控制器设计的线性化状态空间方程。相比第三章的广义动力学形式，该表达式在保留左右轮、左右腿与机体俯仰耦合的同时，把 $s$ 与 $\phi$ 显式纳入状态量，使变量物理含义更直观，也更便于后续求解 LQR 增益并分析控制律。

\section{LQR 主平衡控制器设计}

\subsection{性能指标构造}

对于式 \eqref{eq:linear-state-space} 描述的十维四输入线性时不变系统，本文采用线性二次型调节器（LQR）作为主平衡控制器\cite{klemm2020lqr}。其思路是通过权衡状态偏差与控制输入，求解一组使综合代价最小的状态反馈增益，用于在平衡点附近实现稳定与跟踪。

定义无穷时域二次性能指标为
\begin{equation}
  J =
  \int_{0}^{\infty}
  \left(
    \delta \boldsymbol{x}^{\mathsf{T}}\boldsymbol{Q}\,\delta \boldsymbol{x}
    +
    \delta \boldsymbol{u}^{\mathsf{T}}\boldsymbol{R}\,\delta \boldsymbol{u}
  \right)\dd t,
  \label{eq:lqr-cost}
\end{equation}
其中，$\boldsymbol{Q} \succeq 0$ 为状态加权矩阵，$\boldsymbol{R} \succ 0$ 为控制加权矩阵。$\boldsymbol{Q}$ 用于约束前进位移、偏航角、左右腿摆角与机体俯仰角等状态偏离平衡点的程度，实际取值上通常对 $\dot{s}$、$\phi$ 与 $\theta_b$ 给予较大权重；$\boldsymbol{R}$ 用于限制四个驱动通道的输入幅值，降低饱和与冲击风险。

\subsection{最优状态反馈律推导}

对于系统
\begin{equation}
  \delta \dot{\boldsymbol{x}}
  =
  \boldsymbol{A}\delta \boldsymbol{x}
  +
  \boldsymbol{B}\delta \boldsymbol{u},
  \label{eq:lqr-system}
\end{equation}
若存在正定矩阵 $\boldsymbol{P}$ 满足连续时间代数 Riccati 方程
\begin{equation}
  \boldsymbol{A}^{\mathsf{T}}\boldsymbol{P}
  +
  \boldsymbol{P}\boldsymbol{A}
  -
  \boldsymbol{P}\boldsymbol{B}\boldsymbol{R}^{-1}\boldsymbol{B}^{\mathsf{T}}\boldsymbol{P}
  +
  \boldsymbol{Q}
  =
  \boldsymbol{0},
  \label{eq:riccati}
\end{equation}
则使式 \eqref{eq:lqr-cost} 最小化的最优状态反馈律为
\begin{equation}
  \delta \boldsymbol{u}
  =
  -\boldsymbol{K}\delta \boldsymbol{x},
  \qquad
  \boldsymbol{K} =
  \boldsymbol{R}^{-1}\boldsymbol{B}^{\mathsf{T}}\boldsymbol{P}.
  \label{eq:lqr-law}
\end{equation}

因此，主平衡控制器输出为
\begin{equation}
  \boldsymbol{u}^{\ast}
  =
  \begin{bmatrix}
    T_{lw,l}^{\ast} \\
    T_{lw,r}^{\ast} \\
    T_{bl,l}^{\ast} \\
    T_{bl,r}^{\ast}
  \end{bmatrix}
  =
  -
  \boldsymbol{K}
  \left(
    \boldsymbol{x} - \boldsymbol{x}_d
  \right),
  \label{eq:lqr-output}
\end{equation}
其中，$\boldsymbol{x}_d$ 为期望状态向量。在静态站立任务中，通常取 $\boldsymbol{x}_d = \boldsymbol{0}$；在速度跟踪任务中，则可通过设置 $\dot{s}_d$、$s_d$ 及 $\phi_d$ 等参考量生成对应的平衡状态。

在直线平衡与匀速行进工况下，由于机器人结构左右对称且参考轨迹关于中心面对称，通常进一步引入公共模态量
\begin{equation}
  T_{w}^{\ast}
  =
  \frac{T_{lw,l}^{\ast} + T_{lw,r}^{\ast}}{2},
  \qquad
  T_{\theta}^{\ast}
  =
  \frac{T_{bl,l}^{\ast} + T_{bl,r}^{\ast}}{2},
  \label{eq:lqr-common-mode}
\end{equation}
作为后续矢状面控制分配和腿部任务空间控制的名义输入。

\subsection{LQR 控制量的物理意义}

式 \eqref{eq:lqr-output} 表明，LQR 控制器直接输出左右轮电机力矩和左右腿摆驱动力矩四个分量。其中，$T_{lw,l}^{\ast}$ 与 $T_{lw,r}^{\ast}$ 共同决定整机前进加速度与偏航修正；$T_{bl,l}^{\ast}$ 与 $T_{bl,r}^{\ast}$ 则主要参与左右腿摆姿态调整，并通过改变机体质心相对支撑点的几何关系来辅助俯仰稳定。对于本文关注的直线平衡工况，上述四个输入又可进一步提炼为式 \eqref{eq:lqr-common-mode} 中的公共轮系力矩 $T_w^\ast$ 与公共腿摆力矩 $T_\theta^\ast$，从而与后续腿部 VMC 控制环节自然衔接。

\section{腿部控制器设计}

\subsection{腿长通道控制律}

LQR 主控制器主要处理矢状面内的平衡与前进，而腿部还需要完成支撑、缓冲与高度调节等任务。因此，本文在腿部任务空间引入 VMC 控制器，对腿长与腿摆分别设计控制律。

对于左、右腿分别定义任务空间状态
\begin{equation}
  \boldsymbol{x}_{v,i} =
  \begin{bmatrix}
    l_i & \theta_{l,i}
  \end{bmatrix}^{\mathsf{T}},
  \qquad i \in \{l,r\},
  \label{eq:leg-task-state}
\end{equation}
其中，$l_i$ 表示第 $i$ 条腿的虚拟腿长，$\theta_{l,i}$ 表示第 $i$ 条腿的虚拟腿摆角。首先在腿长方向构造虚拟弹簧阻尼器，其轴向虚拟力为
\begin{equation}
  F_{l,i}^\ast
  =
  K_l (l_{d,i} - l_i)
  +
  D_l (\dot{l}_{d,i} - \dot{l}_i),
  \qquad i \in \{l,r\},
  \label{eq:leg-force-law}
\end{equation}
其中，$l_{d,i}$ 为期望腿长。式 \eqref{eq:leg-force-law} 使腿部在支撑阶段表现为具有等效刚度和阻尼的虚拟支撑单元，从而能够在保持机体高度的同时吸收落地冲击并抑制腿长振荡。

\subsection{腿摆通道控制律}

在腿摆方向，本文将 LQR 主控制器输出的公共摆动控制量 $T_\theta^\ast$ 作为主导项，同时保留局部阻尼校正，以形成腿摆控制律
\begin{equation}
  T_{\theta,i}^\ast
  =
  T_\theta^\ast
  +
  K_{\theta p}(\theta_{d,i} - \theta_{l,i})
  +
  D_{\theta}(\dot{\theta}_{d,i} - \dot{\theta}_{l,i}),
  \qquad i \in \{l,r\},
  \label{eq:leg-swing-law}
\end{equation}
其中，$\theta_{d,i}$ 为局部腿摆参考角。对于直线平衡工况，通常取 $\theta_{d,l} = \theta_{d,r}$；当机器人执行跨障、下蹲或收腿等动作时，则可在该通道叠加动作规划器给出的摆角参考。

将轴向虚拟力和摆动虚拟转矩组合，可得第 $i$ 条腿的任务空间控制量
\begin{equation}
  \boldsymbol{F}_{v,i}^\ast
  =
  \begin{bmatrix}
    F_{l,i}^\ast \\
    T_{\theta,i}^\ast
  \end{bmatrix},
  \qquad i \in \{l,r\}.
  \label{eq:task-space-input}
\end{equation}

\subsection{任务空间到关节空间的力矩映射}

根据第三章推导得到的虚拟腿雅可比矩阵 $\boldsymbol{J}_{v,i}$，由虚功原理可将任务空间控制量映射为关节空间驱动力矩
\begin{equation}
  \boldsymbol{\tau}_i^\ast
  =
  \boldsymbol{J}_{v,i}^{\mathsf{T}}\boldsymbol{F}_{v,i}^\ast,
  \qquad i \in \{l,r\}.
  \label{eq:leg-joint-torque}
\end{equation}
式 \eqref{eq:leg-joint-torque} 给出了腿部控制器在执行层的表达式。该映射建立在五连杆机构运动学关系之上，可将任务空间控制量转换为关节力矩命令，从而驱动实际闭环腿部机构。

\section{滚转补偿器设计}

\subsection{滚转补偿需求分析}

第三章的主动力学模型侧重矢状面内的前进与俯仰平衡，但在实际运动中，地面不平、侧向扰动以及转向时的离心效应都可能引起机体滚转。若滚转不加补偿，左右腿支撑力分配容易失衡，严重时会出现单侧车轮卸载甚至失去稳定接触。为此，本文在主控制器之外增加滚转补偿通道，对左右腿轴向支撑输出做差动修正。

\subsection{虚拟滚转力矩构造}

记机体滚转角为 $\gamma$，滚转角速度为 $\dot{\gamma}$，则可构造期望虚拟滚转补偿力矩
\begin{equation}
  M_{\mathrm{roll}}^\ast
  =
  K_{r}(\gamma_d - \gamma)
  +
  D_{r}(\dot{\gamma}_d - \dot{\gamma}),
  \label{eq:roll-moment}
\end{equation}
其中，$\gamma_d$ 通常取为 0，表示期望机体保持横向水平姿态；$K_r$ 和 $D_r$ 分别为滚转补偿的比例系数和微分系数。

设机器人左右腿支撑点相对机体中心线的半轮距为 $b$，则由静力矩平衡关系可得
\begin{equation}
  M_{\mathrm{roll}}^\ast
  =
  b \left(
    F_{l,r}^{\mathrm{roll}} - F_{l,l}^{\mathrm{roll}}
  \right),
  \label{eq:roll-force-balance}
\end{equation}
其中，$F_{l,l}^{\mathrm{roll}}$ 和 $F_{l,r}^{\mathrm{roll}}$ 分别为左右腿因滚转补偿而引入的附加轴向力。为保持附加力之和为零，从而不改变整机总支撑力，可令
\begin{equation}
  F_{l,l}^{\mathrm{roll}} = -\Delta F_{\mathrm{roll}},\qquad
  F_{l,r}^{\mathrm{roll}} = \Delta F_{\mathrm{roll}},
  \label{eq:roll-force-symmetry}
\end{equation}
联立式 \eqref{eq:roll-force-balance} 和式 \eqref{eq:roll-force-symmetry} 可得
\begin{equation}
  \Delta F_{\mathrm{roll}}
  =
  \frac{M_{\mathrm{roll}}^\ast}{2b}.
  \label{eq:deltaF-roll}
\end{equation}

\subsection{与腿部 VMC 的融合方式}

将式 \eqref{eq:deltaF-roll} 产生的差动轴向力叠加至腿长控制通道，可得滚转补偿后的左右腿轴向控制量
\begin{equation}
  F_{l,l}^\ast = F_l^\ast - \Delta F_{\mathrm{roll}},
  \qquad
  F_{l,r}^\ast = F_l^\ast + \Delta F_{\mathrm{roll}},
  \label{eq:roll-force-fusion}
\end{equation}
其中，$F_l^\ast$ 为腿长主控制器输出的公共轴向支撑力。当机体向左侧倾斜时，滚转补偿器将增大右腿支撑力并减小左腿支撑力，从而在机体横滚轴上形成恢复力矩；反之亦然。由此，滚转补偿器无需改变矢状面 LQR 主控制器结构，仅通过对左右腿 VMC 轴向力分配进行修正，即可实现侧向姿态稳定功能。

\section{控制器综述}

\subsection{控制量融合}

综合上述设计，本文控制器的输出由 LQR 主平衡通道、腿部任务空间通道以及滚转补偿通道共同组成。在忽略偏航差动力矩的直线对称工况下，左右轮控制命令可写为
\begin{equation}
  T_{lw,l}^\ast = T_{lw,r}^\ast = T_w^\ast.
  \label{eq:wheel-command}
\end{equation}

左右腿任务空间目标则分别表示为
\begin{equation}
  \boldsymbol{F}_{v,l}^\ast
  =
  \begin{bmatrix}
    F_l^\ast - \Delta F_{\mathrm{roll}} \\
    T_\theta^\ast
  \end{bmatrix},
  \qquad
  \boldsymbol{F}_{v,r}^\ast
  =
  \begin{bmatrix}
    F_l^\ast + \Delta F_{\mathrm{roll}} \\
    T_\theta^\ast
  \end{bmatrix}.
  \label{eq:left-right-task-command}
\end{equation}
最后经由式 \eqref{eq:leg-joint-torque} 映射为实际关节力矩命令，下发至腿部执行器。

\subsection{控制层级说明}

本文控制器可概括为如下层级结构：
\begin{enumerate}
  \item 状态估计层：融合编码器、惯性测量单元和腿部解算结果，实时获得 $s$、$\phi$、$\theta_{ll}$、$\theta_{lr}$、$\theta_b$、$\gamma$ 及其导数；
  \item 主平衡控制层：基于线性化状态空间模型和 LQR 反馈律，生成 $T_{lw,l}^\ast$、$T_{lw,r}^\ast$、$T_{bl,l}^\ast$ 与 $T_{bl,r}^\ast$，并在直线工况下提取公共模态量 $T_w^\ast$ 和 $T_\theta^\ast$；
  \item 腿部任务空间控制层：基于 VMC 生成公共腿长支撑力 $F_l^\ast$ 以及左右腿任务空间控制量；
  \item 滚转补偿层：根据滚转角与滚转角速度生成差动轴向补偿量 $\Delta F_{\mathrm{roll}}$；
  \item 执行层：将任务空间控制量映射为关节空间力矩，并与轮系力矩命令一起发送至电机驱动器。
\end{enumerate}

上述结构将“矢状面主平衡”和“腿部任务空间调节”分开处理。LQR 主控制器利用平衡点附近的线性化模型，以较低计算量实现姿态稳定与行进控制；腿部 VMC 以任务空间的虚拟力/转矩为核心，便于用少量参数完成支撑与摆动调节；滚转补偿器则用于补足侧向姿态稳定需求。三者组合后，在不显著增加计算负担的前提下，可以同时覆盖平衡、支撑与抗扰控制目标。

\subsection{本章小结}

本章给出了整机控制器的完整设计流程：在第三章动力学模型基础上，对五自由度耦合模型在平衡点附近进行线性化，并结合 $s$、$\phi$ 的运动学映射构造状态空间方程；基于二次型性能指标设计 LQR 主平衡控制器，得到最优状态反馈律；在腿部任务空间内采用 VMC，分别给出腿长与腿摆通道控制律，并通过虚拟腿雅可比矩阵实现任务空间到关节空间的力矩映射；此外，加入基于差动支撑力分配的滚转补偿器以提升侧向稳定性。以上各部分共同构成本文双轮足机器人混合控制方案。
